% ═══════════════════════════════════════════════════════════════
% SPANISCH LEHRERHINWEISE LH-03b: hay + muebles
% ═══════════════════════════════════════════════════════════════
% Sequenz 3: Mi casa | Block b | Klassenstufe 7
% Kurzplan + Differenzierung + Lösungen
% ═══════════════════════════════════════════════════════════════

\documentclass[11pt, a4paper]{article}

% ═══════════════════════════════════════════════════════════════
% SW-MODUS TOGGLE
% ═══════════════════════════════════════════════════════════════
\newif\ifswmodus
\swmodusfalse

% ═══════════════════════════════════════════════════════════════
% PAKETE
% ═══════════════════════════════════════════════════════════════

\usepackage[utf8]{inputenc}
\usepackage[T1]{fontenc}
\usepackage[ngerman]{babel}
\usepackage{helvet}
\renewcommand{\familydefault}{\sfdefault}

\usepackage[margin=2cm]{geometry}
\usepackage{enumitem}
\usepackage{tabularx}
\usepackage{booktabs}
\usepackage{longtable}

\usepackage{xcolor}
\usepackage{amssymb}
\usepackage{tcolorbox}
\tcbuselibrary{skins}

\usepackage{fancyhdr}
\usepackage{lastpage}
\usepackage{needspace}

% ═══════════════════════════════════════════════════════════════
% FARBEN
% ═══════════════════════════════════════════════════════════════

\definecolor{spanisch-color}{HTML}{c41e3a}
\definecolor{grau-color}{HTML}{888888}
\definecolor{hellgrau-color}{HTML}{f5f5f5}
\definecolor{basis-color}{HTML}{2a7a4b}
\definecolor{standard-color}{HTML}{2c5aa0}
\definecolor{erweiterung-color}{HTML}{7b1fa2}
\definecolor{loesung-color}{HTML}{c62828}

\definecolor{sw-dark}{HTML}{000000}
\definecolor{sw-gray}{HTML}{666666}
\definecolor{sw-white}{HTML}{ffffff}

\ifswmodus
    \colorlet{spanisch}{sw-dark}
    \colorlet{grau}{sw-gray}
    \colorlet{hellgrau}{sw-white}
    \colorlet{basis}{sw-dark}
    \colorlet{standard}{sw-dark}
    \colorlet{erweiterung}{sw-dark}
    \colorlet{loesung}{sw-dark}
\else
    \colorlet{spanisch}{spanisch-color}
    \colorlet{grau}{grau-color}
    \colorlet{hellgrau}{hellgrau-color}
    \colorlet{basis}{basis-color}
    \colorlet{standard}{standard-color}
    \colorlet{erweiterung}{erweiterung-color}
    \colorlet{loesung}{loesung-color}
\fi

\newcommand{\fachfarbe}{spanisch}

% ═══════════════════════════════════════════════════════════════
% VARIABLEN
% ═══════════════════════════════════════════════════════════════

\newcommand{\thema}{hay + muebles (Es gibt + Möbel)}
\newcommand{\sequenz}{03}
\newcommand{\block}{b}
\newcommand{\fach}{Spanisch}
\newcommand{\klasse}{7}
\newcommand{\dauer}{90 Min. (Doppelstunde)}
\newcommand{\lerngruppengroesse}{24}
\newcommand{\datum}{\today}

% ═══════════════════════════════════════════════════════════════
% KOPF- UND FUSSZEILE
% ═══════════════════════════════════════════════════════════════

\pagestyle{fancy}
\fancyhf{}
\renewcommand{\headrulewidth}{0.4pt}
\renewcommand{\footrulewidth}{0.4pt}

\lhead{\fach{} -- Klasse \klasse}
\chead{\textbf{LH-\sequenz\block{} (Lehrerhinweise)}}
\rhead{\datum}

\lfoot{}
\cfoot{}
\rfoot{Seite \thepage\ von \pageref{LastPage}}

% ═══════════════════════════════════════════════════════════════
% BEFEHLE
% ═══════════════════════════════════════════════════════════════

\newcommand{\B}{\textcolor{basis}{\textbf{[B]}}}
\newcommand{\Std}{\textcolor{standard}{\textbf{[S]}}}
\newcommand{\E}{\textcolor{erweiterung}{\textbf{[E]}}}

\newcommand{\loesung}[1]{\textcolor{loesung}{\textbf{#1}}}

\newcommand{\es}[1]{\textcolor{\fachfarbe}{\textbf{#1}}}

\newtcolorbox{kurzplan}{
    colback=hellgrau,
    colframe=\fachfarbe,
    colbacktitle=white,
    coltitle=\fachfarbe,
    boxrule=1pt,
    arc=0pt,
    fonttitle=\bfseries,
    attach boxed title to top left={yshift=-2mm, xshift=5mm},
    boxed title style={boxrule=0pt, colback=white},
    title=Kurzplan
}

\newtcolorbox{differenzierung}{
    colback=white,
    colframe=grau,
    colbacktitle=white,
    coltitle=grau,
    boxrule=0.5pt,
    arc=0pt,
    fonttitle=\bfseries,
    attach boxed title to top left={yshift=-2mm, xshift=5mm},
    boxed title style={boxrule=0pt, colback=white},
    title=Differenzierung
}

\newtcolorbox{fehlerbox}{
    colback=white,
    colframe=loesung,
    colbacktitle=white,
    coltitle=loesung,
    boxrule=0.5pt,
    arc=0pt,
    fonttitle=\bfseries,
    attach boxed title to top left={yshift=-2mm, xshift=5mm},
    boxed title style={boxrule=0pt, colback=white},
    title=Häufige Fehler
}

\newtcolorbox{materialbox}{
    colback=white,
    colframe=\fachfarbe,
    colbacktitle=white,
    coltitle=\fachfarbe,
    boxrule=0.5pt,
    arc=0pt,
    fonttitle=\bfseries,
    attach boxed title to top left={yshift=-2mm, xshift=5mm},
    boxed title style={boxrule=0pt, colback=white},
    title=Material-Checkliste
}

% ═══════════════════════════════════════════════════════════════
% DOKUMENT
% ═══════════════════════════════════════════════════════════════

\begin{document}

% ═══════════════════════════════════════════════════════════════
% KOPFBEREICH
% ═══════════════════════════════════════════════════════════════

\begin{center}
    {\Large\bfseries\textcolor{\fachfarbe}{Lehrerhinweise: \thema}}\\[0.3em]
    {\normalsize LH-\sequenz\block{} | \dauer}
\end{center}

\vspace{10pt}

% ═══════════════════════════════════════════════════════════════
% 1. KURZPLAN
% ═══════════════════════════════════════════════════════════════

\section*{1. Stundenverlauf (Kurzplan)}

\textbf{1. Stunde (45 Min.) -- Einführung Wortschatz + hay}

\begin{tabularx}{\textwidth}{|c|l|X|l|}
\hline
\textbf{Zeit} & \textbf{Phase} & \textbf{Inhalt} & \textbf{Material} \\
\hline
0--5 & Einstieg & Bildimpuls: eingerichtetes Zimmer zeigen, Frage: ``¿Qué hay?'' & Bild/Beamer \\
\hline
5--15 & Input & 8 Möbelwörter mit Bildkarten einführen, chorisches Sprechen & Bildkarten \\
\hline
15--25 & Übung 1 & Einzelabfrage: Karte zeigen $\rightarrow$ ``¿Qué es?'' $\rightarrow$ ``Es una/un...'' & Bildkarten \\
\hline
25--35 & Erarbeitung & Struktur \es{hay + un/una} einführen, Beispielsätze an Tafel & Tafel \\
\hline
35--45 & Übung 2 & PA: ``¿Qué hay en tu dormitorio?'' -- Antworten sammeln & -- \\
\hline
\end{tabularx}

\vspace{1em}

\textbf{2. Stunde (45 Min.) -- Anwendung + Produktion}

\begin{tabularx}{\textwidth}{|c|l|X|l|}
\hline
\textbf{Zeit} & \textbf{Phase} & \textbf{Inhalt} & \textbf{Material} \\
\hline
0--5 & Wiederholung & Schnelles Bildquiz: Karten hochhalten, SuS rufen Wort & Bildkarten \\
\hline
5--15 & Vertiefung & AB erklären, Fokus: ``Nach hay kommt un/una!'' & AB \\
\hline
15--30 & Übung & EA: AB Aufgaben 1-3 bearbeiten, L geht herum & AB \\
\hline
30--40 & Produktion & EA: Aufgabe 4 (eigenes Zimmer beschreiben) & AB \\
\hline
40--45 & Abschluss & 2-3 SuS lesen vor, freundliche Korrektur, HA ansagen & -- \\
\hline
\end{tabularx}

% ═══════════════════════════════════════════════════════════════
% 2. DIFFERENZIERUNG
% ═══════════════════════════════════════════════════════════════

\section*{2. Differenzierung}

\begin{differenzierung}
\textbf{Legende:} \B{} Basis (BR) | \Std{} Standard (MR) | \E{} Erweiterung (GY)

\vspace{0.5em}

\textbf{WICHTIG:} Diese Marker sind NUR für die Lehrkraft. Auf Schülermaterial erscheinen KEINE Niveaumarkierungen!
\end{differenzierung}

\vspace{1em}

\begin{tabularx}{\textwidth}{|l|X|X|X|}
\hline
& \B{} \textbf{Basis} & \Std{} \textbf{Standard} & \E{} \textbf{Erweiterung} \\
\hline
\textbf{Aufgabe 1} & Mit Bildkarten als Hilfe & Selbstständig & + Sätze mit hay bilden \\
\hline
\textbf{Aufgabe 2} & Wortliste sichtbar lassen & Selbstständig & + Plural (hay dos sillas) \\
\hline
\textbf{Aufgabe 3} & 1 Satz pro Zimmer reicht & Vollständige Sätze & + Adjektive (grande, nuevo) \\
\hline
\textbf{Aufgabe 4} & 2-3 Sätze, Vorlage nutzen & Mind. 4 Sätze & + 6 Sätze, auch ``No hay...'' \\
\hline
\end{tabularx}

\vspace{1em}

\textbf{Konkrete Hinweise:}
\begin{itemize}
    \item \B{} SuS mit LRS: Zeitzugabe (+10 Min.), Merkkasten darf offen bleiben
    \item \B{} DaZ-SuS: Möbelwörter vor der Stunde auf Deutsch-Spanisch besprechen
    \item \E{} Leistungsstarke: Zusatzaufgabe ``Beschreibe das Klassenzimmer!'' oder weitere Möbel recherchieren
    \item \E{} Schnelle: Als Experten einsetzen, helfen bei PA
\end{itemize}

% ═══════════════════════════════════════════════════════════════
% 3. HÄUFIGE FEHLER
% ═══════════════════════════════════════════════════════════════

\section*{3. Häufige Fehler}

\begin{fehlerbox}
\begin{tabularx}{\textwidth}{|l|X|X|}
\hline
\textbf{Fehler} & \textbf{Beispiel} & \textbf{Reaktion} \\
\hline
Bestimmter Artikel nach hay & *``Hay \underline{la} cama'' statt \es{hay una cama} & Kontrastieren: ``hay = es gibt (neu)'' vs. ``La cama está = Das Bett ist (bekannt)'' \\
\hline
Falsches Genus (un/una) & *``hay \underline{un} mesa'' statt \es{hay una mesa} & Artikel-Regel wiederholen: la $\rightarrow$ una, el $\rightarrow$ un \\
\hline
Betonung sofá & *``SÓfa'' statt \es{soFÁ} & Übertrieben betonen, Akzent zeigen, chorisches Sprechen \\
\hline
Schreibung lámpara & *``lampara'' ohne Akzent & Auf Dreisilbigkeit hinweisen: LÁM-pa-ra \\
\hline
hay/está verwechselt & *``Hay en mi dormitorio'' & Regel: hay + Objekt (Existenz), estar + Ort (Position) \\
\hline
\end{tabularx}
\end{fehlerbox}

% ═══════════════════════════════════════════════════════════════
% 4. MATERIAL-CHECKLISTE
% ═══════════════════════════════════════════════════════════════

\section*{4. Material-Checkliste}

\begin{materialbox}
\begin{itemize}[label=$\square$]
    \item AB-\sequenz\block.pdf (\lerngruppengroesse{} Kopien)
    \item ML-\sequenz\block.pdf (1$\times$ für Lehrkraft)
    \item Lehrwerk: Qué pasa 1, S. 38-42
    \item Bildkarten Möbel (8 Stück, A5-Format): mesa, silla, sofá, cama, armario, lámpara, estantería, escritorio
    \item Bild/Foto eines eingerichteten Zimmers (für Einstieg)
    \item Tafelmarker / Kreide
    \item Optional: LP-\sequenz\block.html (Tablets/PCs für Übungsphase oder HA)
    \item Optional: Magnete für Bildkarten an Tafel
\end{itemize}
\end{materialbox}

% ═══════════════════════════════════════════════════════════════
% 5. LÖSUNGEN
% ═══════════════════════════════════════════════════════════════

\section*{5. Lösungen AB-\sequenz\block}

\textbf{Merkkasten 1: hay}
\begin{itemize}
    \item hay bedeutet: \loesung{es gibt}
    \item En mi dormitorio hay una cama = \loesung{In meinem Schlafzimmer gibt es ein Bett.}
    \item En el salón hay un sofá = \loesung{Im Wohnzimmer gibt es ein Sofa.}
\end{itemize}

\textbf{Merkkasten 2: Los muebles}

\begin{tabular}{ll|ll}
la mesa & = \loesung{der Tisch} & el armario & = \loesung{der Schrank} \\
la silla & = \loesung{der Stuhl} & la lámpara & = \loesung{die Lampe} \\
el sofá & = \loesung{das Sofa} & la estantería & = \loesung{das Regal} \\
la cama & = \loesung{das Bett} & el escritorio & = \loesung{der Schreibtisch} \\
\end{tabular}

\vspace{1em}

\textbf{Aufgabe 1: Zuordnung} (4 Punkte)

\begin{tabular}{ll}
\es{la cama} & $\rightarrow$ \loesung{das Bett} \\
\es{el escritorio} & $\rightarrow$ \loesung{der Schreibtisch} \\
\es{la estantería} & $\rightarrow$ \loesung{das Regal} \\
\es{el armario} & $\rightarrow$ \loesung{der Schrank} \\
\end{tabular}

\vspace{1em}

\textbf{Aufgabe 2: hay + un/una} (4 Punkte)

\begin{tabular}{ll}
a) & En mi dormitorio \loesung{hay una} cama grande. \\
b) & En el salón \loesung{hay un} sofá nuevo. \\
c) & En la cocina \loesung{hay una} mesa. \\
d) & En mi habitación \loesung{hay un} escritorio. \\
\end{tabular}

\vspace{1em}

\textbf{Aufgabe 3: Zimmer beschreiben} (6 Punkte)

\begin{tabular}{ll}
a) Wohnzimmer: & \loesung{En el salón hay un sofá y una lámpara.} \\
b) Küche: & \loesung{En la cocina hay una mesa y cuatro sillas.} \\
c) Arbeitszimmer: & \loesung{En el despacho hay un escritorio, una estantería y una lámpara.} \\
\end{tabular}

\textit{Auch akzeptabel: ``En la oficina...'', ``En mi habitación...'', andere Formulierungen mit korrekter Struktur.}

\vspace{1em}

\textbf{Aufgabe 4: Eigenes Zimmer} (6 Punkte)

Individuelle Lösungen. Bewertungskriterien:
\begin{itemize}
    \item Mindestens 4 Sätze mit hay: 2P
    \item Korrekter Artikel (un/una): 2P
    \item Korrekte Möbelvokabeln: 1P
    \item Sprachliche Korrektheit: 1P
\end{itemize}

% ═══════════════════════════════════════════════════════════════
% 6. HAUSAUFGABE
% ═══════════════════════════════════════════════════════════════

\section*{6. Hausaufgabe}

\begin{itemize}
    \item Lehrwerk S. 40, Aufgabe 3 (Sätze mit hay bilden)
    \item Vokabeln lernen: 8 Möbelwörter mit Artikel
    \item Optional: Lernpfad LP-\sequenz\block.html bearbeiten
    \item \E{} Zusatz: 5 weitere Möbelwörter recherchieren (z.B. el espejo, la alfombra, la cortina)
\end{itemize}

% ═══════════════════════════════════════════════════════════════
% 7. REFLEXIONSNOTIZEN
% ═══════════════════════════════════════════════════════════════

\section*{7. Reflexionsnotizen (nach Durchführung)}

\begin{tabularx}{\textwidth}{|l|X|}
\hline
\textbf{Was lief gut?} & \\[2em]
\hline
\textbf{Was lief nicht?} & \\[2em]
\hline
\textbf{Zeitmanagement} & \\[2em]
\hline
\textbf{Für nächstes Mal} & \\[2em]
\hline
\end{tabularx}

\vspace{1em}

\textbf{Spezifische Fragen:}
\begin{itemize}
    \item Haben alle SuS die Regel ``hay + un/una'' verstanden?
    \item Welche Möbel bereiteten Ausspracheschwierigkeiten?
    \item War Aufgabe 4 in der Zeit schaffbar?
    \item Wurde die Differenzierung genutzt?
\end{itemize}

\end{document}
